\scanpage{172}
% archive sheet number: 171
\deleted{\unsure{En te}} Si $\underline G$ est total dans $\underline U$, \inserted{i.e. $A$ est une algèbre de fonctions sur $G$,} alors la classe $\underline K^{U}$ s'identifie,
avec \deleted{toutes} ses structures, à la classe des \deleted{\unsure{rep}}
$G$-modules \inserted{de dim finie} tels que les \unsure{coéfficients|coëfficients|coefficients} \deleted{\unsure{en sont|sont}}
\deleted{\unsure{membres}} soient $\in A$.

Ainsi, une $\underline K$-hyperalgèbre $\underline U$ définit
un \uline{groupe} $G$, et une algèbre \inserted{\unsure{avec $1$|avec $1$,}} $A_0(G)$ de fonctions
sur $G$ à valeurs dans $k$, séparant les $\mathrm{p}^{\mathrm{ts}}$ de $G$,
\margin{\deleted{\unsure{à vérifier}}}
telle que \unsure{tt|ts} caractère \deleted{\unsure{que tt caractère de $A_0$ soit}} de $A_0$ soit défini
\deleted{stable sous $\nu$, et \unsure{ses translatées à g. et à droite}}
\inserted{par un $\mathrm{p}^{\mathrm{t}}$ de $G$, $A_0$ stable par $\unsure{\nu|v|\vee}$, et telle que}
% slanted
\margin{ceci \uline{implique} que $A_0$ est stable sous translations g. et dr.}
telle que $f\in A_0$ implique $f(xy)\in A_0\otimes A_0$. Alors
tout $\underline U$-module permis \uline{définit} un $\underline G$
module, permis pour $A_0$, [et la réciproque
est vraie si $G$ est \deleted{dense} \inserted{total} dans $U$].

\uline{Renversons la vapeur} partons d'un groupe
$G$ avec une algèbre $A_0$ ayant les propriétés
précédentes. Considérons la catégorie des
$G$-modules permis. Il lui correspond
une $\underline K$-algèbre $\underline U$, et
\diagram{\begin{tikzcd}[ampersand replacement=\&,row sep=small]
G \arrow[r] \& \underline U.\\
A_0 \& \underline A \arrow[l]
\end{tikzcd}}

Je dis que $G\to\underline U$ est une bijection
de $G$ sur \deleted{\unsure{un}} $G(\underline U)$ (éléments inversibles
de $U$ tels que $\delta(s)=s\otimes s$). Que ce
soit injectif résulte du fait que $A_0$
sépare les $\mathrm p^{\mathrm{ts}}$ de $G$, il est trivial que $G$
est appliqué dans $G(\underline U)$, \deleted{\unsure{p|pour}} montrons que
c'est surjectif. Comme tt caractère de $A_0$
provient de $G$, il suffit de prouver que $A\to A_0$
